\section{隔离易感者回流}
\label{sec:sup:return}
正文研究隔离易感者在清零前不释放的情形（第~\ref{sec:model}~节）。本节用数值说明：若被隔离的未感染者在疫情过程中陆续返回社区，正文的哪些结论仍成立，哪些会改变。

\subsection{模型与设定}
含回流的模型在式\eqref{eq:model:full}的 $S$、$S_q$ 方程中加入回流项 $\omega S_q$：
\begin{equation}
\begin{cases}
\dot S=-\dfrac{c[\beta+(1-\beta)q]}{N}SI+\omega S_q,\\[.5em]
\dot I=\dfrac{\beta c(1-q)}{N}SI-\gamma I,\\[.5em]
\dot S_q=\dfrac{(1-\beta)cq}{N}SI-\omega S_q,\\[.5em]
\dot I_q=\dfrac{\beta cq}{N}SI-\delta_q I_q,\\[.5em]
\dot R=\gamma I+\delta_q I_q,
\end{cases}
\qquad\omega\ge0.
\label{eq:return:model}
\end{equation}
回流率 $\omega$ 为常数，隔离易感者的停留时间服从均值 $1/\omega$ 的指数分布，$\omega=0$ 即式\eqref{eq:model:full}。本节取 $1/\omega=14$ d，与既有隔离时长研究涉及的时间尺度相近\cite{Ashcroft2021,Wells2021}。指数停留是简化的回流机制，与固定期限或检测后解除不同。

阈值策略的实施方式与正文相同。$I$ 首次到达 $\eta$ 时启动；此后取 $c_0S_c/S\le c\le c_0$，隔离比例由平衡式\eqref{eq:model:balance}确定，从而维持 $I\equiv\eta$；$S$ 首次降到 $S_c$ 时恢复 $(c_0,q_0)$ 并一直保持。退出时不额外释放隔离易感者，$S_q$ 继续按回流机制演化，退出后也不重新启动干预。有回流时阈值控制期内的 $S$ 可能回升，上述接触率范围对 $S\ge S_c$ 照样适用。两组参数的初值同正文，$S_q(0)=0$。累计新增感染计至退出后 $I$ 首次降到 $1$ 的时刻
\begin{equation}
T_{\rm cut}=\inf\{t>t_2:I(t)=1,\ \dot I(t)<0\},
\label{eq:return:cut}
\end{equation}
与正文清零时间的约定相同。

追踪与回流只在 $S$ 与 $S_q$ 之间转移人数。两式相加，总易感者 $U=S+S_q$ 满足 $\dot U=-\beta cSI/N$，因此
\begin{equation}
\Itcum(T)=U(0)-U(T).
\label{eq:return:U}
\end{equation}

比较五种控制：仅隔离控制、仅减少接触、无回流成本最优律（正文定理~\ref{thm:joint:cost-minimum}的状态反馈，$w_c=1$、$w_q=2$，直接用于回流模型），以及在回流模型中重新优化的控制时长最短与成本 $J$ 最低的控制。重新优化在阈值策略内进行，用时间域 Radau 配点，区间数由 $20$ 逐级加密到 $160$，粗网格取十个起点。最终候选用完整五仓室方程回放核验，并用分段常控制的打靶法复核。

\subsection{数值结果}
回流下控制时长最短的控制收敛到仅隔离控制，表~\ref{tab:sup:return:term}中不再单列。
\begin{table}[htbp]
\centering\small
\renewcommand{\arraystretch}{1.2}
\begin{threeparttable}
\caption{平均停留 $14$ d 的指数回流情景下各控制的结果}
\label{tab:sup:return:term}
\begin{tabular*}{0.8\linewidth}{@{\extracolsep{\fill}}lccc@{}}
\toprule
控制 & $\Delta t$（d） & $J$（d） & $\Itcum(T_{\rm cut})$\\
\midrule
\multicolumn{4}{@{}l}{基准参数（人）}\\
仅隔离控制 & $\nrRwBaseQDt$ & $\nrRwBaseQJ$ & $\nrRwBaseQInf$\\
回流下的成本最优控制 & $\nrRwBaseODt$ & $\nrRwBaseOJ$ & $\nrRwBaseOInf$\\
无回流成本最优律 & $\nrRwBaseMDt$ & $\nrRwBaseMJ$ & $\nrRwBaseMInf$\\
仅减少接触 & $\nrRwBaseCDt$ & $\nrRwBaseCJ$ & $\nrRwBaseCInf$\\
\midrule
\multicolumn{4}{@{}l}{西安参数（$\eta=0.001N$，$10^6$ 人）}\\
仅隔离控制 & $\nrRwXianQDt$ & $\nrRwXianQJ$ & $\nrRwXianQInf$\\
回流下的成本最优控制 & $\nrRwXianODt$ & $\nrRwXianOJ$ & $\nrRwXianOInf$\\
无回流成本最优律 & $\nrRwXianMDt$ & $\nrRwXianMJ$ & $\nrRwXianMInf$\\
仅减少接触 & $\nrRwXianCDt$ & $\nrRwXianCJ$ & $\nrRwXianCInf$\\
\bottomrule
\end{tabular*}
\begin{tablenotes}[flushleft]\footnotesize
\item 注：$1/\omega=14$ d；$\Itcum(T_{\rm cut})$ 见式\eqref{eq:return:cut}；$J$ 取 $w_c=1$、$w_q=2$。
\end{tablenotes}
\end{threeparttable}
\end{table}

正文的结构性结论在这一情景中仍成立：
\begin{itemize}
\item 仅隔离控制的控制时长最短，仅减少接触最长；
\item 回流下的成本最优控制先仅隔离（基准约前 $\nrRwBaseFrontPct\%$、西安约前 $\nrRwXianFrontPct\%$ 的控制期），随后两种措施并用；直接沿用无回流成本最优律，成本只高约 $\nrRwBaseMigGapPct\%$ 与 $\nrRwXianMigGapPct\%$；
\item 各控制退出后 $I$ 都未越过 $\eta$。
\end{itemize}

改变的是量级。仅隔离控制的控制时长由无回流的 $\nrRwBaseQDtNoret$ d 与 $\nrRwXianQDtNoret$ d 增至 $\nrRwBaseQDt$ d 与 $\nrRwXianQDt$ d。累计新增感染接近 $U(0)-S_c$，后者基准为 $\nrRwBaseL$ 人、西安为 $\nrRwXianL\times10^6$ 人。控制之间的感染差别也明显缩小：仅减少接触比仅隔离控制多出的感染，基准参数下由无回流时的约 $\nrRwBaseDiffPctNoret\%$ 降到约 $\nrRwBaseDiffPct\%$，西安参数下由约 $\nrRwXianDiffPctNoret\%$ 降到约 $\nrRwXianDiffPct\%$。

原因在于终止时总易感者已低于常规传播临界值：各轨迹在 $T_{\rm cut}$ 时 $U/S_c$ 在基准参数下为 $\nrRwBaseUMin$--$\nrRwBaseUMax$，西安参数下约为 $\nrRwXianUMax$。由式\eqref{eq:return:U}，累计新增感染因而都超过 $U(0)-S_c$。无回流时，大量未感染接触者一直留在 $S_q$，退出时 $S$ 降到 $S_c$ 并不需要这么多感染；有回流时他们返回社区后仍会被感染，追踪隔离对他们的保护是暂时的。
